% DGHNOUTS - mathematical documentation of the spectral/hp library
%
% Notation and equation numbers follow
%   G. Karniadakis and S. Sherwin, Spectral/hp Element Methods for
%   Computational Fluid Dynamics, 2nd ed., Oxford University Press, 2005.
% Equations taken from the book carry the book's number, e.g. (A.1.3).
% Equations that are not in the book are numbered (D.1), (D.2), ...
% The standard-element coordinate is xi, as in the book's Chapter 2
% (Appendix A writes x).
%
% Build:  latexmk -pdf dghnouts_theory.tex     (or: pdflatex twice)
\documentclass[11pt,a4paper]{article}

\usepackage[utf8]{inputenc}
\usepackage[margin=2.5cm]{geometry}
\usepackage{amsmath,amssymb}
\usepackage{booktabs}
\usepackage{array}
\usepackage{xcolor}
\usepackage{listings}
\usepackage[colorlinks=true,linkcolor=blue!50!black,citecolor=blue!50!black,urlcolor=blue!50!black]{hyperref}

% ---------------------------------------------------------------------------
% Notation (Karniadakis & Sherwin)
% ---------------------------------------------------------------------------
\newcommand{\Jac}[3]{P_{#1}^{#2,#3}}            % P_n^{alpha,beta}
\newcommand{\Pab}[1]{P_{#1}^{\alpha,\beta}}     % P_n^{alpha,beta}
\newcommand{\dd}{\mathrm{d}}
\newcommand{\ddxi}{\frac{\dd}{\dd\xi}}
\newcommand{\code}[1]{\texttt{#1}}

% equations from the book keep the book's numbers; ours are (D.n)
\newcounter{dgeq}
\newcommand{\dgtag}{\refstepcounter{dgeq}\tag{D.\arabic{dgeq}}}

\lstset{
  language=C++, basicstyle=\ttfamily\small, columns=fullflexible,
  keywordstyle=\color{blue!60!black}\bfseries, commentstyle=\color{green!40!black},
  frame=single, framerule=0.3pt, rulecolor=\color{black!30},
  xleftmargin=1em, xrightmargin=1em, breaklines=true,
}

\title{\textbf{DGHNOUTS}\\[0.3em]
  \large Mathematical documentation of the spectral/$hp$ library\\[0.2em]
  \normalsize Part~I: Jacobi polynomials}
\author{Carlos Junqueira Junior\\ \small DynFluid, Arts et M\'etiers -- ENSAM}
\date{\today}

\begin{document}
\maketitle

\begin{abstract}
This document gives the mathematical definition of every function of the
DGHNOUTS spectral/$hp$ library, the equation each one implements, and the
identities used to test it. Notation and equation numbers follow Karniadakis
and Sherwin~\cite{KS2005}: an equation taken from the book keeps its number,
for example~(A.1.3); an equation that is not in the book, derived here, is
numbered (D.1), (D.2), \dots\ The document grows with the code; this part
covers the module \code{include/jacobi.hpp}.
\end{abstract}

\tableofcontents

% ===========================================================================
\section{Notation and conventions}
% ===========================================================================

\begin{itemize}
  \item $\xi \in [-1,1]$ is the coordinate of the standard interval, as in the
        book's Chapter~2. Appendix~A of the book writes the same polynomials
        in a generic variable $x$; here $\xi$ is used throughout, and $x$ is
        kept for the physical coordinate, $x=\chi(\xi)$, which will appear
        with the element mappings.
  \item $\Pab{n}(\xi)$ is the Jacobi polynomial of degree $n \ge 0$ with
        weight parameters $\alpha,\beta > -1$.
  \item $L_n(\xi) = \Jac{n}{0}{0}(\xi)$ is the Legendre polynomial and
        $T_n(\xi)$ the Chebyshev polynomial.
  \item $\Gamma(z)$ is the Gamma function, $\Gamma(n+1)=n!$ for integer $n\ge0$.
\end{itemize}

\paragraph{Names in the code.}
\begin{center}
\begin{tabular}{@{}lll@{}}
\toprule
Symbol & Code & Meaning \\
\midrule
$\xi$, $\xi_i$ & \code{xi}, \code{xi[i]} & standard-element coordinate, points \\
$n$ & \code{n} & degree (target degree of a routine) \\
$\alpha$, $\beta$ & \code{alpha}, \code{beta} & Jacobi weight parameters \\
$n_p$ & \code{np} & number of points \\
$a_n^1,\dots,a_n^4$ & \code{a1}, \dots, \code{a4} & coefficients of (A.1.3) \\
$\xi_1,\xi_2,\xi_3$ & \code{xi1}, \code{xi2}, \code{xi3} & standard-hexahedron coordinates (future) \\
\bottomrule
\end{tabular}
\end{center}
In three dimensions the book writes the hexahedron coordinates
$\xi_1,\xi_2,\xi_3$ and reserves $\eta_1,\eta_2,\eta_3$ for the collapsed
coordinates of triangles, tetrahedra and prisms. DGHNOUTS follows the same
convention, so \code{eta} never denotes a second hexahedron direction.

\paragraph{Code conventions.}
All functions live in namespace \code{dg}. They throw
\code{std::invalid\_argument} when $n<0$ or when $\alpha\le-1$ or
$\beta\le-1$, since the polynomials are only defined (orthogonal) for
$\alpha,\beta>-1$. Array routines take a raw pointer and a size,
\code{(const T* xi, std::size\_t np)}, so they accept any contiguous storage
(\code{std::vector}, C arrays, Kokkos host views).

\paragraph{Floating-point precision.}
\label{sec:precision}
Every function is a template on the floating-point type \code{T} of the points
and results. \code{T} is deduced from \code{xi}; for the norm, which has no
point argument, it is given explicitly:
\begin{center}
\code{jacobi\_norm2<float>(n, alpha, beta)}.
\end{center}
The library is compiled for \code{float}, \code{double} and \code{long double}. The
parameters $\alpha,\beta$ are always \code{double}; the recurrence
coefficients are computed in \code{T}. Table~\ref{tab:precision} gives the
largest relative errors measured against \code{long double} by test~9
(Section~\ref{sec:tests}).

\begin{table}[h]
\centering
\begin{tabular}{@{}lcc@{}}
\toprule
& \code{float} & \code{double} \\
\midrule
values $\Pab{n}(\xi_i)$                    & $5\times10^{-6}$ & $2\times10^{-14}$ \\
derivatives $\bigl[\Pab{n}\bigr]'(\xi_i)$  & $7\times10^{-5}$ & $8\times10^{-14}$ \\
norms $h_n^{\alpha,\beta}$                 & $1\times10^{-5}$ & $2\times10^{-14}$ \\
\midrule
unit round-off                              & $6\times10^{-8}$ & $1.1\times10^{-16}$ \\
\bottomrule
\end{tabular}
\caption{Largest relative errors against \code{long double}, for
$n=0,\dots,20$, all parameter pairs of Section~\ref{sec:tests} and 43 points
including $\xi=\pm1$ (x86, 80-bit \code{long double}).}
\label{tab:precision}
\end{table}

The recurrence amplifies the rounding error of each operation by a factor of
a few hundred at $n=20$, and the norm suffers from cancellation between large
$\ln\Gamma$ terms, so in \code{float} only 4 to 5 digits are correct. Nodes,
weights and matrices for a single-precision solver should therefore be
computed in \code{double} and rounded to \code{float} afterwards, which keeps
them accurate to $6\times10^{-8}$. The \code{long double} type is 80-bit
extended precision on x86 (unit round-off $5\times10^{-20}$) and IEEE
quadruple precision on Arm processors such as Rhea or Grace
(software-emulated, slower); it can be used to compute high-order nodes and
weights to full \code{double} accuracy.

% ===========================================================================
\section{Jacobi polynomials}
\label{sec:jacobi}
% ===========================================================================

\subsection{Definition}

Jacobi polynomials $\Pab{n}(\xi)$ are the polynomial solutions of the singular
Sturm--Liouville problem~\cite[App.~A.1]{KS2005}
\begin{equation}
  (1-\xi)(1+\xi)\,\frac{\dd^2 y}{\dd\xi^2}
  + \bigl[\beta-\alpha-(\alpha+\beta+2)\xi\bigr]\frac{\dd y}{\dd\xi}
  = -\lambda_n\, y,
  \tag{A.1.1}
\end{equation}
with
\begin{equation*}
  \lambda_n = n(n+\alpha+\beta+1),
  \qquad y(\xi) = \Pab{n}(\xi).
\end{equation*}
They can be written explicitly with the Rodrigues formula
\begin{equation*}
  \Pab{n}(\xi) = \frac{(-1)^n}{2^n n!}\,(1-\xi)^{-\alpha}(1+\xi)^{-\beta}
  \frac{\dd^n}{\dd\xi^n}\Bigl[(1-\xi)^{\alpha+n}(1+\xi)^{\beta+n}\Bigr],
  \qquad \alpha,\beta>-1,
\end{equation*}
but this form is never used for computation; the recurrence of
Section~\ref{sec:values} is.

\paragraph{Special cases.}
\begin{align*}
  \alpha=\beta=0 &: \quad L_n(\xi) = \Jac{n}{0}{0}(\xi)
     && \text{(Legendre)},\\
  \alpha=\beta=-\tfrac12 &: \quad T_n(\xi) = \frac{2^{2n}(n!)^2}{(2n)!}\,
     \Jac{n}{-1/2}{-1/2}(\xi) && \text{(Chebyshev)}.
\end{align*}
In a DG code the values that matter are $\alpha=\beta=0$ (Legendre modal basis
and Gauss--Legendre points), $\alpha=\beta=1$ (interior Gauss--Lobatto points,
Section~\ref{sec:deriv}), and $(\alpha,\beta)=(2i+1,0)$ or $(2i+2j+2,0)$ for
the collapsed-coordinate bases of simplices.

% ===========================================================================
\subsection{Values: the three-term recurrence}
\label{sec:values}
% ===========================================================================

The polynomials are evaluated with the recursion relation
\begin{equation}
\begin{aligned}
  \Jac{0}{\alpha}{\beta}(\xi) &= 1,\\
  \Jac{1}{\alpha}{\beta}(\xi) &= \tfrac12\bigl[\alpha-\beta+(\alpha+\beta+2)\xi\bigr],\\
  a_n^1\,\Pab{n+1}(\xi) &= \bigl(a_n^2 + a_n^3 \xi\bigr)\Pab{n}(\xi) - a_n^4\,\Pab{n-1}(\xi),
\end{aligned}
\tag{A.1.3}
\end{equation}
with
\begin{align*}
  a_n^1 &= 2(n+1)(n+\alpha+\beta+1)(2n+\alpha+\beta),\\
  a_n^2 &= (2n+\alpha+\beta+1)(\alpha^2-\beta^2),\\
  a_n^3 &= (2n+\alpha+\beta)(2n+\alpha+\beta+1)(2n+\alpha+\beta+2),\\
  a_n^4 &= 2(n+\alpha)(n+\beta)(2n+\alpha+\beta+2).
\end{align*}
For $n\ge1$ and $\alpha,\beta>-1$ we have $2n+\alpha+\beta>0$ and
$n+\alpha+\beta+1>0$, so $a_n^1\neq0$ and the recursion never divides by zero.

\paragraph{Implementation.}
The internal helper \code{rec\_coef(n, alpha, beta)} returns
$a_n^1,\dots,a_n^4$ exactly as in~(A.1.3). Inside a routine, however,
\code{n} is the target degree, so the recurrence step is indexed by
\code{k}: step $k$ is (A.1.3) with the book's $n$ replaced by $k$. It calls
\code{rec\_coef(k, alpha, beta)} and produces $\Pab{k+1}$ from $\Pab{k}$ and
$\Pab{k-1}$, for $k=1,\dots,n-1$.

For $n_p$ points, the array routine runs the degree loop outside and the point
loop inside:
\begin{lstlisting}
for (int k = 1; k < n; ++k) {
    const RecCoef<T> c = rec_coef(k, a, b);       // once per k
    for (std::size_t i = 0; i < np; ++i) {        // independent: vectorised
        const T pkp1 = ((c.a2 + c.a3*xi[i])*p[i] - c.a4*pm1[i]) / c.a1;
        pm1[i] = p[i];  p[i] = pkp1;
    }
}
\end{lstlisting}
Each point still needs its own recurrence, so the operation count is
$O(n\,n_p)$ as with $n_p$ scalar calls; the gain comes from SIMD
vectorisation of the inner loop and from computing the coefficients once per
degree instead of once per point.

\begin{center}
\begin{tabular}{@{}ll@{}}
\toprule
Function & Computes \\
\midrule
\code{T jacobi(n, alpha, beta, xi)} & $\Pab{n}(\xi)$ at one point \\
\code{void jacobi(n, alpha, beta, xi, np, p)} & $p_i=\Pab{n}(\xi_i)$, $i=0,\dots,n_p-1$ \\
\code{void jacobi(n, alpha, beta, xi, p)} & same, \code{std::vector} overload \\
\bottomrule
\end{tabular}
\end{center}

% ===========================================================================
\subsection{Derivatives}
\label{sec:deriv}
% ===========================================================================

\paragraph{The book's derivative relation, and why it is not used.}
Appendix~A gives the derivative through
\begin{equation}
\begin{aligned}
  b_n^1(\xi)\,\ddxi\Pab{n}(\xi) &= b_n^2(\xi)\,\Pab{n}(\xi) + b_n^3(\xi)\,\Pab{n-1}(\xi),\\
  b_n^1(\xi) &= (2n+\alpha+\beta)(1-\xi^2),\\
  b_n^2(\xi) &= n\bigl[\alpha-\beta-(2n+\alpha+\beta)\xi\bigr],\\
  b_n^3(\xi) &= 2(n+\alpha)(n+\beta).
\end{aligned}
\tag{A.1.4}
\end{equation}
This is what Polylib's \code{jacobfd} uses. Since $b_n^1(\pm1)=0$, it cannot
give the derivative at the end points $\xi=\pm1$, which are nodes of every
Gauss--Lobatto rule. DGHNOUTS does not use~(A.1.4).

\paragraph{Derivative by degree reduction.}
Instead, \code{jacobi\_deriv} uses
\begin{equation}
  \ddxi\Pab{n}(\xi) = \frac12\,(\alpha+\beta+n+1)\,\Jac{n-1}{\alpha+1}{\beta+1}(\xi),
  \tag{A.1.8}
\end{equation}
with $\ddxi\Jac{0}{\alpha}{\beta}=0$. The right-hand side is evaluated with the
recurrence~(A.1.3) for the parameters $(\alpha+1,\beta+1)$, so it is valid on
the whole closed interval $[-1,1]$. Applying~(A.1.8) twice gives the second
derivative used in the tests:
\begin{equation}
  \frac{\dd^2}{\dd\xi^2}\Pab{n}(\xi)
  = \frac14\,(n+\alpha+\beta+1)(n+\alpha+\beta+2)\,
    \Jac{n-2}{\alpha+2}{\beta+2}(\xi).
  \dgtag\label{eq:second-deriv}
\end{equation}

Relation~(A.1.8) also explains the Gauss--Lobatto points of Appendix~B: the
interior points are the zeros of $\bigl[\Jac{Q-1}{\alpha}{\beta}\bigr]'$, that
is, of $\Jac{Q-2}{\alpha+1}{\beta+1}$. For Legendre, they are the zeros of
$\Jac{Q-2}{1}{1}$.

\paragraph{Values and derivatives in one pass.}
Differentiating~(A.1.3) with respect to $\xi$ (the coefficients $a_n^j$ do not
depend on $\xi$) gives a recurrence for the derivative that runs in the same
loop as the values:
\begin{equation}
  a_n^1\,\bigl[\Pab{n+1}\bigr]'(\xi)
  = a_n^3\,\Pab{n}(\xi)
  + \bigl(a_n^2 + a_n^3 \xi\bigr)\bigl[\Pab{n}\bigr]'(\xi)
  - a_n^4\,\bigl[\Pab{n-1}\bigr]'(\xi),
  \dgtag\label{eq:deriv-rec}
\end{equation}
started with
\begin{equation}
  \bigl[\Jac{0}{\alpha}{\beta}\bigr]' = 0,
  \qquad
  \bigl[\Jac{1}{\alpha}{\beta}\bigr]' = \tfrac12(\alpha+\beta+2).
  \dgtag
\end{equation}
Like~(A.1.8), relation~\eqref{eq:deriv-rec} contains no division by
$1-\xi^2$ and holds at $\xi=\pm1$. It is what \code{jacobi\_and\_deriv} uses.
In exact arithmetic it gives the same derivative as~(A.1.8); in double
precision the two agree to about $10^{-14}$ relative error for $n\le20$.

\paragraph{Implementation.}
\code{jacobi\_and\_deriv} keeps four work arrays
($\Pab{k}$, $\Pab{k-1}$ and their derivatives). For large $n_p$ these arrays
would be streamed through main memory at every step $k$, which makes the
routine slower than two separate calls. The points are therefore processed in
blocks of 256, so that the four arrays stay in the L1 cache for all $n$ steps.

\begin{center}
\begin{tabular}{@{}lll@{}}
\toprule
Function & Computes & Equation \\
\midrule
\code{T jacobi\_deriv(n, alpha, beta, xi)} & $\bigl[\Pab{n}\bigr]'(\xi)$ & (A.1.8) \\
\code{void jacobi\_deriv(n, alpha, beta, xi, np, dp)} & $\bigl[\Pab{n}\bigr]'(\xi_i)$ & (A.1.8) \\
\code{void jacobi\_and\_deriv(n, alpha, beta,} & $\Pab{n}(\xi_i)$ and & (A.1.3), \eqref{eq:deriv-rec} \\
\quad\code{xi, np, p, dp)} & $\bigl[\Pab{n}\bigr]'(\xi_i)$ & \\
\bottomrule
\end{tabular}
\end{center}
In \code{jacobi\_and\_deriv}, either output pointer may be \code{nullptr}
when only one of the two results is needed.

% ===========================================================================
\subsection{All degrees at once: the generalised Vandermonde matrix}
\label{sec:vandermonde}
% ===========================================================================

Each step of~(A.1.3) produces the next degree, so a single pass up to
$N=$\code{nmax} gives every polynomial $\Pab{0},\dots,\Pab{N}$. Evaluated at
points $\xi_0,\dots,\xi_{n_p-1}$, they form the generalised Vandermonde matrix
\begin{equation}
  \mathbf{V}_{ik} = \Pab{k}(\xi_i),
  \qquad i=0,\dots,n_p-1, \quad k=0,\dots,N.
  \dgtag\label{eq:vandermonde}
\end{equation}
With Legendre polynomials ($\alpha=\beta=0$) and $\xi_i$ the nodes of a DG
element, $\mathbf{V}$ maps modal coefficients $\hat u_k$ to nodal values,
$u(\xi_i)=\sum_k \mathbf{V}_{ik}\hat u_k$, which is the transform needed for
modal filtering and shock sensors.

\paragraph{Implementation.}
\code{jacobi\_all(nmax, alpha, beta, xi, np, V)} stores $\mathbf{V}$ in
row-major order, one row per point and one column per degree:
\begin{equation*}
  \code{V[i*(nmax+1) + k]} = \mathbf{V}_{ik} = \Pab{k}(\xi_i).
\end{equation*}
The array \code{V} must hold $n_p(N+1)$ values. Each $\Pab{k}$ is computed
exactly once, whereas calling \code{jacobi} for each degree
$k=0,\dots,N$ would repeat the first $k$ steps of the recurrence every time.

% ===========================================================================
\subsection{Orthogonality and norm}
\label{sec:norm}
% ===========================================================================

Jacobi polynomials are orthogonal on $[-1,1]$ with respect to the weight
$(1-\xi)^\alpha(1+\xi)^\beta$:
\begin{equation}
\begin{gathered}
  \int_{-1}^{1}(1-\xi)^\alpha(1+\xi)^\beta\,\Pab{n}(\xi)\,\Pab{m}(\xi)\,\dd\xi = 0,
  \qquad n\neq m,\\[0.4em]
  \int_{-1}^{1}(1-\xi)^\alpha(1+\xi)^\beta\,\Pab{n}(\xi)\,\Pab{n}(\xi)\,\dd\xi
  = \frac{2^{\alpha+\beta+1}}{2n+\alpha+\beta+1}\,
    \frac{\Gamma(n+\alpha+1)\,\Gamma(n+\beta+1)}{n!\,\Gamma(n+\alpha+\beta+1)}
  \equiv h_n^{\alpha,\beta}.
\end{gathered}
\tag{A.1.7}
\end{equation}
\code{jacobi\_norm2(n, alpha, beta)} returns $h_n^{\alpha,\beta}$.

\paragraph{The case $n=0$.}
For $n=0$ and $\alpha+\beta=-1$ (e.g.\ Chebyshev), formula~(A.1.7) reads
$0\cdot\infty$: the factor $2n+\alpha+\beta+1$ vanishes and
$\Gamma(n+\alpha+\beta+1)=\Gamma(0)$ is infinite. Using
$z\,\Gamma(z)=\Gamma(z+1)$ with $z=n+\alpha+\beta+1$,
\begin{equation*}
  (2n+\alpha+\beta+1)\,\Gamma(n+\alpha+\beta+1)
  \;\xrightarrow{\;n=0\;}\;
  (\alpha+\beta+1)\,\Gamma(\alpha+\beta+1) = \Gamma(\alpha+\beta+2),
\end{equation*}
so that
\begin{equation}
  h_0^{\alpha,\beta}
  = \frac{2^{\alpha+\beta+1}\,\Gamma(\alpha+1)\,\Gamma(\beta+1)}{\Gamma(\alpha+\beta+2)},
  \dgtag\label{eq:h0}
\end{equation}
which is finite for all $\alpha,\beta>-1$. The code uses~\eqref{eq:h0} for
$n=0$ and~(A.1.7) for $n\ge1$.

\paragraph{Overflow.}
$\Gamma(z)$ overflows double precision for $z>171.6$. The norm is therefore
computed from logarithms with \code{std::lgamma} ($\ln|\Gamma|$, C++11),
\begin{equation*}
\begin{split}
  \ln h_n^{\alpha,\beta}
  ={}& (\alpha+\beta+1)\ln2 - \ln(2n+\alpha+\beta+1)
  + \ln\Gamma(n+\alpha+1) + \ln\Gamma(n+\beta+1)\\
  & - \ln\Gamma(n+1) - \ln\Gamma(n+\alpha+\beta+1),
\end{split}
\end{equation*}
and $h_n^{\alpha,\beta}=\exp(\ln h_n^{\alpha,\beta})$. The terms
$\ln\Gamma$ are large and nearly cancel, which is why the norm loses more
accuracy in \code{float} than the values do (Table~\ref{tab:precision}).

\paragraph{Checks.}
For Legendre, (A.1.7) reduces to $h_n^{0,0}=2/(2n+1)$. For Chebyshev,
\eqref{eq:h0} gives $h_0^{-1/2,-1/2}=\Gamma(\tfrac12)^2/\Gamma(1)=\pi$, the
classical value of $\int_{-1}^{1}(1-\xi^2)^{-1/2}\,\dd\xi$.

% ===========================================================================
\section{Verification}
\label{sec:tests}
% ===========================================================================

The program \code{tests/test\_jacobi.cpp} checks every function against
identities of Appendix~A, for degrees $n=0,\dots,20$ and for the parameter
pairs
\begin{equation*}
  (\alpha,\beta)\in\bigl\{(0,0),\,(1,0),\,(2,0),\,(1,1),\,(-\tfrac12,-\tfrac12),\,
  (\tfrac12,\tfrac32),\,(0.3,-0.7)\bigr\}.
\end{equation*}
The last pair, with non-integer parameters, is beyond what Polylib can handle.
Tests 1 to 8 run in \code{double}; test 9 compares the three precisions. The
points are rounded to \code{float}, so that they are exactly the same in every
precision.

\paragraph{End-point value.}
\begin{equation}
  \Pab{n}(1) = \binom{n+\alpha}{n} = \frac{(n+\alpha)!}{\alpha!\,n!},
  \tag{A.1.5}
\end{equation}
evaluated for non-integer $\alpha$ as
$\Gamma(n+\alpha+1)/\bigl(\Gamma(\alpha+1)\,n!\bigr)$.

\paragraph{Symmetry.}
\begin{equation}
  \Pab{n}(-\xi) = (-1)^n\,\Jac{n}{\beta}{\alpha}(\xi).
  \tag{A.1.6}
\end{equation}

\paragraph{Differential equation.}
The residual of~(A.1.1) is computed with $y'$ from~(A.1.8) and $y''$
from~\eqref{eq:second-deriv}, at interior points and at $\xi=\pm1$.

\paragraph{End-point derivative.}
Combining (A.1.8) and (A.1.5) with $\alpha=\beta=0$ gives
\begin{equation}
  L_n'(1) = \tfrac12(n+1)\,\Jac{n-1}{1}{1}(1) = \tfrac12(n+1)\,n,
  \dgtag
\end{equation}
which checks that the derivative is correct exactly where~(A.1.4) fails.

\paragraph{Consistency of the array routines.}
\code{jacobi} and \code{jacobi\_deriv} on arrays, and every column of
\code{jacobi\_all}, must agree with the single-point routines to the last bit
(they perform the same operations). \code{jacobi\_and\_deriv} must agree to
$10^{-12}$, since its derivative comes from~\eqref{eq:deriv-rec} rather
than~(A.1.8), and must give bitwise the same results when one of its outputs
is \code{nullptr}.

\paragraph{Precision.}
The values, derivatives and norms in \code{float} and \code{double} are
compared with \code{long double} (Table~\ref{tab:precision}). The tolerances
are set from the measured errors: $10^{-5}$ (values) and $10^{-4}$
(derivatives, norms) in \code{float}, $10^{-13}$ to $10^{-12}$ in
\code{double}. The \code{long double} instantiation is checked against the
end-point value~(A.1.5) to $10^{-16}$, below what \code{double} can reach,
which confirms that it really computes in extended precision.

\paragraph{Optional comparison with Polylib.}
When CMake is given \code{-DPOLYLIB\_DIR=\dots}, the tests also compare values
and interior derivatives with Polylib's \code{jacobfd}, and
check~(A.1.7) numerically with Polylib's Gauss--Jacobi quadrature
\code{zwgj}, which integrates exactly polynomials of degree $\le 2Q-1$.

\begin{center}
\begin{tabular}{@{}llll@{}}
\toprule
Test & Identity & Functions checked & Tolerance \\
\midrule
1 & Polylib \code{jacobfd} values & \code{jacobi} & $10^{-13}$ \\
2 & Polylib \code{jacobfd} derivatives & \code{jacobi\_deriv} & $10^{-11}$ \\
3 & (A.1.5) & \code{jacobi} & $10^{-13}$ \\
4 & (A.1.6) & \code{jacobi} & $10^{-13}$ \\
5 & (A.1.7) & \code{jacobi}, \code{jacobi\_norm2} & $10^{-12}$ \\
6 & (A.1.1) & \code{jacobi}, \code{jacobi\_deriv} & $10^{-11}$ \\
7 & $L_n'(1)=n(n+1)/2$ & \code{jacobi\_deriv} & $10^{-13}$ \\
8 & array vs.\ single point, \code{nullptr} & all array routines & $0$ / $10^{-12}$ \\
9 & \code{float}, \code{double} vs.\ \code{long double} & all routines & Table~\ref{tab:precision} \\
\bottomrule
\end{tabular}
\end{center}
Tests 1, 2 and 5 require Polylib. Errors are relative, with a floor of $1$
so that values close to zero are compared absolutely.

% ===========================================================================
\appendix
\section{Function reference}
% ===========================================================================

All functions are declared as \code{template <typename T>}, with
\code{T} $\in\{$\code{float}, \code{double}, \code{long double}$\}$.

\begin{center}
\small
\begin{tabular}{@{}>{\raggedright\arraybackslash}p{0.47\textwidth}>{\raggedright\arraybackslash}p{0.30\textwidth}l@{}}
\toprule
Signature & Result & Equation \\
\midrule
\code{T jacobi(int n, double alpha, double beta, T xi)}
  & $\Pab{n}(\xi)$ & (A.1.3) \\
\code{T jacobi\_deriv(int n, double alpha, double beta, T xi)}
  & $\bigl[\Pab{n}\bigr]'(\xi)$ & (A.1.8) \\
\code{void jacobi(int n, double alpha, double beta, const T* xi, std::size\_t np, T* p)}
  & $p_i=\Pab{n}(\xi_i)$ & (A.1.3) \\
\code{void jacobi\_deriv(int n, double alpha, double beta, const T* xi, std::size\_t np, T* dp)}
  & $dp_i=\bigl[\Pab{n}\bigr]'(\xi_i)$ & (A.1.8) \\
\code{void jacobi\_and\_deriv(int n, double alpha, double beta, const T* xi, std::size\_t np, T* p, T* dp)}
  & $p_i$ and $dp_i$ in one pass & (A.1.3), \eqref{eq:deriv-rec} \\
\code{void jacobi\_all(int nmax, double alpha, double beta, const T* xi, std::size\_t np, T* V)}
  & $\mathbf{V}_{ik}=\Pab{k}(\xi_i)$, row-major & \eqref{eq:vandermonde} \\
\code{T jacobi\_norm2<T = double>(int n, double alpha, double beta)}
  & $h_n^{\alpha,\beta}$ & (A.1.7), \eqref{eq:h0} \\
\bottomrule
\end{tabular}
\end{center}
The \code{std::vector<T>} overloads of \code{jacobi} and \code{jacobi\_deriv}
resize the output to the size of \code{xi} and call the pointer versions.

% ===========================================================================
\begin{thebibliography}{9}
\bibitem{KS2005}
  G.~Karniadakis and S.~Sherwin,
  \emph{Spectral/hp Element Methods for Computational Fluid Dynamics},
  2nd ed., Oxford University Press, 2005.
\bibitem{AS1964}
  M.~Abramowitz and I.~A.~Stegun,
  \emph{Handbook of Mathematical Functions}, Chapter~22,
  National Bureau of Standards, 1964.
\end{thebibliography}

\end{document}
